% ============================ INTRODUCTION ============================
\chapter*{Introduction générale}
\addcontentsline{toc}{chapter}{Introduction générale}
\markboth{INTRODUCTION GÉNÉRALE}{}

La transformation numérique du secteur bancaire impose aux établissements de
disposer de systèmes d'information capables, d'une part, de centraliser la
gestion de leur réseau d'agences et, d'autre part, d'exploiter leurs données
pour piloter l'activité. La maîtrise de l'information --- effectifs, portefeuille
clients, objectifs commerciaux --- constitue aujourd'hui un levier déterminant de
performance et d'aide à la décision.

C'est dans ce contexte que s'inscrit ce projet de fin d'études, réalisé au sein
de la \textbf{BTK -- Banque Tuniso-Koweïtienne} en vue de l'obtention de la
Licence en Business Computing, parcours Business Intelligence, à l'Esprit School
of Business. Le projet consiste à concevoir et réaliser une application web de
gestion des agences bancaires, adossée à un volet décisionnel d'aide à la
décision.

L'application couvre la gestion des agences, des employés, des clients et des
objectifs commerciaux, le \textbf{pointage (présence) des employés}, ainsi qu'un
ensemble de circuits de validation (workflows) encadrant la création et la
modification des données sensibles. À ces fonctionnalités transactionnelles
s'ajoute un volet décisionnel : une chaîne \textbf{ETL} consolide les données dans
un datamart qui alimente des rapports Microsoft Power BI, un tableau de bord
embarqué et une \textbf{segmentation des agences} par apprentissage automatique.

Ce rapport, qui rend compte d'un projet conduit selon la méthodologie
\textbf{CRISP-DM}, est organisé en cinq chapitres.
Le \textbf{premier chapitre} présente le cadre général du projet : l'organisme
d'accueil, l'étude de l'existant, la problématique et la méthodologie adoptée. Le
\textbf{deuxième chapitre} est consacré à l'analyse et à la spécification des
besoins (acteurs, besoins fonctionnels et non fonctionnels, cas d'utilisation).
Le \textbf{troisième chapitre} détaille la conception : architecture, modèle de
données et scénarios dynamiques. Le \textbf{quatrième chapitre} est dédié au
\textbf{volet décisionnel} : la chaîne ETL, le modèle en étoile, les tableaux de
bord et les rapports Power BI, ainsi que la segmentation des agences par
apprentissage automatique. Le \textbf{cinquième et dernier chapitre} présente la
\textbf{réalisation et le déploiement} de la solution. Une conclusion générale
dresse le bilan du travail et ouvre des perspectives d'évolution.
